% Generated draft: requires mathematical and source review.
\documentclass[11pt]{article}
\usepackage[margin=25mm]{geometry}
\usepackage{fontspec}
\setmainfont{Latin Modern Roman}
\usepackage{amsmath,amssymb,mathtools,unicode-math}
\setmathfont{Latin Modern Math}
\usepackage{enumitem,needspace,xurl,hyperref}
\setlength{\parindent}{0pt}
\setlength{\parskip}{5pt}
\setlength{\emergencystretch}{3em}
\newcommand{\N}{\mathbb{N}}
\newcommand{\Z}{\mathbb{Z}}
\newcommand{\Q}{\mathbb{Q}}
\newcommand{\R}{\mathbb{R}}
\newcommand{\C}{\mathbb{C}}
\newcommand{\F}{\mathbb{F}}
\newcommand{\A}{\mathbb{A}}
\newcommand{\PP}{\mathbb{P}}
\newcommand{\E}{\mathbb{E}}
\newcommand{\eps}{\varepsilon}
\newcommand{\dd}{\,\mathrm{d}}
\begin{document}
\section*{UnsolvedMath ID: 30005854}
\textbf{Source problem code:} catalogue entry
\subsection{Definitions and mathematical statement}
Write \((M_-,g_-)\) and \((M_+,g_+)\) schematically for the positive-scalar-curvature manifolds with corners occurring in the referenced Main Theorem. Let \(F_-\subset M_-\) and \(F_+\subset M_+\) denote their distinguished faces, and let \(\phi:F_-\to F_+\) be the prescribed isometry of their induced metrics, so that \(\phi^*(g_+|_{TF_+})=g_-|_{TF_-}\). Attaching the faces identifies \(x\in F_-\) with \(\phi(x)\in F_+\), giving the underlying quotient
\[
M_-\cup_\phi M_+=(M_-\sqcup M_+)/\bigl(x\sim\phi(x)\bigr).
\]
This quotient notation alone specifies neither a smooth structure nor a metric with the required curvature properties.

Let \(H_-\) and \(H_+\) stand schematically for the two face mean-curvature functions mentioned in the source. A function \(H\) does not change sign if \(H\geq0\) everywhere on its face or \(H\leq0\) everywhere on its face. The source does not specify the normal convention or unambiguously identify the faces in its printed mean-curvature subscripts. Let \(\theta_-\) and \(\theta_+\) denote the dihedral-angle functions at corresponding corner points wherever the face identification pairs such points; their conventions and precise domains must be recovered from the referenced theorem.

The supplied problem asks the following three questions, relative to the hypotheses and conclusions of that Main Theorem.
\begin{enumerate}
\item Can its assumption that at least one of \(H_-\) and \(H_+\) does not change sign be removed, retaining its other hypotheses and conclusions?
\item Can an analogue of that theorem be established with the angle hypothesis replaced by the pointwise requirement
\[
\theta_-(p)+\theta_+(\phi(p))=\pi
\]
at every corresponding corner point? Which other hypotheses are retained is not specified in the supplied passage.
\item What sharp hypotheses allow attachment of distinguished isometric faces that need not be cylindrical, while preserving the scalar-curvature and boundary-curvature requirements referred to in the source?
\end{enumerate}
The supplied record omits the Main Theorem and the preceding geometric requirements. Consequently, the dimension restrictions, regularity assumptions, meaning of cylindrical faces, boundary-curvature inequalities, and precise metric or smoothing conclusion cannot be stated from this evidence. These questions therefore remain a conditional problem formulation, rather than a fully specified gluing conjecture.
\subsection{Short English statement}
For positive-scalar-curvature manifolds with corners, ask whether the referenced face-gluing theorem can dispense with its mean-curvature sign assumption, whether pointwise dihedral-angle sums equal to \(\pi\) suffice for an analogue, and which sharp assumptions permit noncylindrical distinguished isometric faces. The omitted theorem and geometric requirements prevent a complete standalone formulation.
\subsection{Sources}
\begin{itemize}
\item Supplied local UnsolvedMath registry record, permanent ID \(30005854\), title \emph{Sharper Positive-Scalar-Curvature Face Gluing}, exact source problem code \texttt{catalogue entry}, in the Oberwolfach Reports Open Problems collection. The record assigns the year 2024; no individual proposer is supplied.
\item \emph{Analysis, Geometry and Topology of Positive Scalar Curvature Metrics} (2024), Oberwolfach Reports, pp.\ 483--566, DOI \texttt{10.4171/owr/2024/9}: \url{https://doi.org/10.4171/owr/2024/9}. This is the supplied report attribution; the extracted passage refers to a Main Theorem and preceding geometric requirements that are absent from this record.
\item \emph{A polyhedron comparison theorem}: \url{https://arxiv.org/abs/1710.08067}. Supplied literature-review reference.
\item \emph{Attaching faces of positive scalar curvature manifolds with corners}: \url{https://arxiv.org/abs/2308.05409}. Supplied literature-review reference.
\item Literature assessment dated exactly \texttt{2026-08-22}: the supplied curation labels the problem partially solved and reports that dihedral-angle comparison, corner-attaching, and polyhedral rigidity results remove or weaken some hypotheses in special settings, while the full sharp sign-free and angle-sum formulation remains open. This is supplied dated curation, not an independently verified literature assessment; its time-sensitive status and cited claims require independent verification.
\item The registry labels the statement \texttt{corrected\_verified} and says that it was reconstructed and checked against the official Oberwolfach report. This is a supplied verification claim, not verification performed for this draft.
\end{itemize}
\end{document}
